% HANDBUCH-SPRECHPRUEFUNGEN-MV.tex
% Bilingual: Deutsch / English (grau)
% Stand: Januar 2026

\documentclass[11pt,a4paper]{article}

% ============================================================================
% PAKETE
% ============================================================================
\usepackage[utf8]{inputenc}
\usepackage[T1]{fontenc}
\usepackage[ngerman]{babel}
\usepackage{helvet}
\renewcommand{\familydefault}{\sfdefault}

\usepackage[margin=2.5cm]{geometry}
\usepackage{xcolor}
\usepackage{tabularx}
\usepackage{booktabs}
\usepackage{longtable}
\usepackage{colortbl}
\usepackage{enumitem}
\usepackage{tikz}
\usepackage{tcolorbox}
\usepackage{hyperref}
\usepackage{fancyhdr}
\usepackage{titlesec}
\usepackage{parskip}
\usepackage{multicol}
\usepackage{amssymb}
\usepackage{pifont}

% ============================================================================
% FARBEN
% ============================================================================
\definecolor{spanischrot}{HTML}{c41e3a}
\definecolor{englishgray}{HTML}{888888}
\definecolor{lightgray}{HTML}{f5f5f5}
\definecolor{bordergray}{HTML}{dddddd}
\definecolor{linkblue}{HTML}{2c5aa0}

% ============================================================================
% HYPERREF SETUP
% ============================================================================
\hypersetup{
    colorlinks=true,
    linkcolor=spanischrot,
    urlcolor=linkblue,
    pdftitle={Handbuch Sprechprüfungen Spanisch MV},
    pdfauthor={}
}

% ============================================================================
% BILINGUALES MAKRO
% ============================================================================
% \bilingual{Deutscher Text}{English text}
\newcommand{\bilingual}[2]{%
    #1\\[2pt]%
    {\small\color{englishgray}\textit{#2}}%
}

% \bipar{Deutscher Absatz}{English paragraph}
\newcommand{\bipar}[2]{%
    #1\par\vspace{2pt}%
    {\small\color{englishgray}\textit{#2}}\par\vspace{8pt}%
}

% \bisection{Deutsche Überschrift}{English heading}
\newcommand{\bisection}[2]{%
    \section{#1}%
    {\small\color{englishgray}\textit{#2}}\par\vspace{6pt}%
}

% \bisubsection{Deutsche Überschrift}{English heading}
\newcommand{\bisubsection}[2]{%
    \subsection{#1}%
    {\small\color{englishgray}\textit{#2}}\par\vspace{4pt}%
}

% \bisubsubsection{Deutsche Überschrift}{English heading}
\newcommand{\bisubsubsection}[2]{%
    \subsubsection{#1}%
    {\small\color{englishgray}\textit{#2}}\par\vspace{2pt}%
}

% ============================================================================
% BOXEN
% ============================================================================
\tcbset{
    colback=lightgray,
    colframe=bordergray,
    boxrule=1pt,
    arc=2pt,
    left=8pt,
    right=8pt,
    top=6pt,
    bottom=6pt
}

\newtcolorbox{infobox}[1][]{
    colback=lightgray,
    colframe=spanischrot,
    title={#1},
    fonttitle=\bfseries,
    coltitle=white,
    colbacktitle=spanischrot
}

\newtcolorbox{warnbox}{
    colback=white,
    colframe=spanischrot,
    boxrule=2pt
}

\newtcolorbox{kritischbox}{
    colback=spanischrot!10,
    colframe=spanischrot,
    boxrule=3pt,
    title={\textbf{RECHTLICH VERBINDLICH / \textit{LEGALLY MANDATORY}}},
    fonttitle=\bfseries\large,
    coltitle=white,
    colbacktitle=spanischrot
}

% ============================================================================
% KOPF- UND FUSSZEILE
% ============================================================================
\pagestyle{fancy}
\fancyhf{}
\fancyhead[L]{\small\color{gray}Handbuch Sprechprüfungen Spanisch MV}
\fancyhead[R]{\small\color{gray}v1.0 | Januar 2026}
\fancyfoot[C]{\thepage}
\renewcommand{\headrulewidth}{0.5pt}

% ============================================================================
% ÜBERSCHRIFTEN
% ============================================================================
\titleformat{\section}
    {\Large\bfseries\color{spanischrot}}
    {\thesection}{1em}{}
\titleformat{\subsection}
    {\large\bfseries}
    {\thesubsection}{1em}{}
\titleformat{\subsubsection}
    {\normalsize\bfseries}
    {\thesubsubsection}{1em}{}

% ============================================================================
% CHECKMARK
% ============================================================================
\newcommand{\cmark}{\ding{51}}
\newcommand{\xmark}{\ding{55}}

% ============================================================================
% DOKUMENT
% ============================================================================
\begin{document}

% ----------------------------------------------------------------------------
% TITELSEITE
% ----------------------------------------------------------------------------
\begin{titlepage}
\centering
\vspace*{2cm}

{\Huge\bfseries\color{spanischrot} Handbuch Sprechprüfungen}\\[0.5cm]
{\Large\color{englishgray}\textit{Oral Examination Handbook}}\\[1.5cm]

{\LARGE Spanisch Mecklenburg-Vorpommern}\\[0.3cm]
{\large\color{englishgray}\textit{Spanish -- Mecklenburg-Western Pomerania}}\\[2cm]

\begin{tcolorbox}[width=12cm, colback=white, colframe=spanischrot]
\centering
{\large\bfseries Komplexe Leistungsermittlung Sprechen}\\
{\color{englishgray}\textit{Complex Assessment of Speaking Skills}}\\[0.5cm]
\&\\[0.5cm]
{\large\bfseries Abitur-Paarprüfung}\\
{\color{englishgray}\textit{Abitur Pair Examination}}
\end{tcolorbox}

\vfill

{\large Version 1.0 | Stand: Januar 2026}\\[0.3cm]
{\color{englishgray}\textit{Version 1.0 | Date: January 2026}}\\[1cm]

{\small Basierend auf den Handreichungen des Ministeriums für Bildung und Kindertagesförderung MV}\\
{\small\color{englishgray}\textit{Based on guidelines from the Ministry of Education MV}}

\end{titlepage}

% ----------------------------------------------------------------------------
% INHALTSVERZEICHNIS
% ----------------------------------------------------------------------------
\tableofcontents
\newpage

% ============================================================================
% EINLEITUNG
% ============================================================================
\bisection{Übersicht: Zwei Prüfungsformate}{Overview: Two Examination Formats}

\bipar{In Mecklenburg-Vorpommern gibt es \textbf{zwei unterschiedliche mündliche Prüfungsformate} für Spanisch in der gymnasialen Oberstufe. Diese werden häufig verwechselt, haben aber verschiedene rechtliche Grundlagen, Zeitpunkte und Funktionen.}{In Mecklenburg-Western Pomerania, there are \textbf{two different oral examination formats} for Spanish in upper secondary education. These are often confused but have different legal bases, timing, and functions.}

\begin{center}
\begin{tikzpicture}
\node[draw=spanischrot, line width=2pt, rounded corners, fill=white, minimum width=14cm, minimum height=6cm] (main) {};

% Titel
\node[anchor=north, font=\bfseries\large] at ([yshift=-0.3cm]main.north) {Mündliche Prüfungen Spanisch MV};
\node[anchor=north, font=\small\itshape, text=englishgray] at ([yshift=-0.8cm]main.north) {Oral Examinations Spanish MV};

% Linke Box
\node[draw=spanischrot, fill=lightgray, rounded corners, minimum width=5.5cm, minimum height=3.5cm, anchor=north west] (left) at ([xshift=0.5cm, yshift=-1.5cm]main.north west) {};
\node[anchor=north, font=\bfseries, text=spanischrot] at ([yshift=-0.2cm]left.north) {Komplexe Leistungs-};
\node[anchor=north, font=\bfseries, text=spanischrot] at ([yshift=-0.6cm]left.north) {ermittlung (KLE)};
\node[anchor=north, font=\small\itshape, text=englishgray] at ([yshift=-1.0cm]left.north) {Complex Assessment};

\node[anchor=north west, font=\footnotesize, align=left] at ([xshift=0.2cm, yshift=-1.5cm]left.north west) {
$\bullet$ Während Q-Phase (Q1--Q3)\\
$\bullet$ Ersetzt eine Klausur\\
$\bullet$ \textbf{PFLICHT} (mind. 1x)\\
$\bullet$ 20 Minuten Dauer\\
$\bullet$ Seit 2019/20
};

% Rechte Box
\node[draw=spanischrot, fill=lightgray, rounded corners, minimum width=5.5cm, minimum height=3.5cm, anchor=north east] (right) at ([xshift=-0.5cm, yshift=-1.5cm]main.north east) {};
\node[anchor=north, font=\bfseries, text=spanischrot] at ([yshift=-0.2cm]right.north) {Abitur-Paarprüfung};
\node[anchor=north, font=\small\itshape, text=englishgray] at ([yshift=-0.6cm]right.north) {Abitur Pair Exam};

\node[anchor=north west, font=\footnotesize, align=left] at ([xshift=0.2cm, yshift=-1.5cm]right.north west) {
$\bullet$ Im Rahmen des Abiturs\\
$\bullet$ Mündl. Prüfungsteil / 4./5. PF\\
$\bullet$ \textbf{OPTIONAL} (auf Antrag)\\
$\bullet$ max. 30 Minuten Dauer\\
$\bullet$ Seit 2025/26
};

\end{tikzpicture}
\end{center}

% ============================================================================
% TEIL A: KOMPLEXE LEISTUNGSERMITTLUNG
% ============================================================================
\newpage
\part*{\color{spanischrot}Teil A: Komplexe Leistungsermittlung Sprechen (KLE)}
\addcontentsline{toc}{part}{Teil A: Komplexe Leistungsermittlung Sprechen}
{\large\color{englishgray}\textit{Part A: Complex Assessment of Speaking Skills}}

\bisection{Rechtliche Grundlagen}{Legal Framework}

\bipar{Die Komplexe Leistungsermittlung Sprechen basiert auf folgenden Rechtsvorschriften:}{The Complex Assessment of Speaking Skills is based on the following regulations:}

\begin{center}
\begin{tabularx}{\textwidth}{|l|X|}
\hline
\rowcolor{spanischrot}\color{white}\textbf{Paragraph} & \color{white}\textbf{Regelung / \textit{Regulation}} \\
\hline
\textbf{§ 22 Abs. 5 APVO} & \textbf{PFLICHT:} Für jede/n SuS erfolgt mindestens eine KLE Sprechen in der gesamten Qualifikationsphase \newline {\color{englishgray}\textit{MANDATORY: At least one KLE Speaking must be conducted for each student during the entire qualification phase}} \\
\hline
\textbf{§ 22 Abs. 4 APVO} & \textbf{ERSATZ:} Max. 1/3 der Klausuren können durch komplexe Leistungen ersetzt werden \newline {\color{englishgray}\textit{REPLACEMENT: Max. 1/3 of written exams may be replaced by complex assessments}} \\
\hline
\textbf{§ 17 Abs. 3 APVO} & \textbf{BEWERTUNG:} Komplexe Leistungen werden wie Klausuren bewertet \newline {\color{englishgray}\textit{GRADING: Complex assessments are graded like written exams}} \\
\hline
\textbf{§ 22 Abs. 2 APVO} & \textbf{AUSNAHME:} Die abiturähnliche Klausur (Q4) kann NICHT ersetzt werden \newline {\color{englishgray}\textit{EXCEPTION: The Abitur-style exam (Q4) CANNOT be replaced}} \\
\hline
\end{tabularx}
\end{center}

\vspace{0.5cm}
\bipar{\textbf{Gültig seit:} 01.08.2020 (Handreichung, VV vom 24.07.2020)}{\textbf{Valid since:} August 1, 2020 (Guidelines, Administrative Regulation dated July 24, 2020)}

\bisubsection{Wann wird die KLE durchgeführt?}{When is the KLE conducted?}

\begin{warnbox}
\centering
\textbf{WICHTIG / \textit{IMPORTANT}}\\[0.3cm]
\begin{tabular}{|c|c|c|c|}
\hline
\rowcolor{lightgray}\textbf{Q1 (11/I)} & \textbf{Q2 (11/II)} & \textbf{Q3 (12/I)} & \textbf{Q4 (12/II)} \\
\hline
Klausur ODER KLE & Klausur ODER KLE & Klausur ODER KLE & \cellcolor{spanischrot!20}\textbf{NUR Klausur!} \\
\hline
\end{tabular}\\[0.3cm]
$\Rightarrow$ \textbf{PFLICHT:} Mind. 1x KLE in Q1, Q2 oder Q3\\
{\color{englishgray}\textit{MANDATORY: At least one KLE in Q1, Q2, or Q3}}\\[0.2cm]
$\Rightarrow$ \textbf{Q4:} Abiturähnliche Klausur -- \textbf{NICHT ersetzbar!}\\
{\color{englishgray}\textit{Q4: Abitur-style exam -- CANNOT be replaced!}}
\end{warnbox}

\bipar{\textbf{Empfehlung:} Die KLE im 2. Schulhalbjahr für 2./3. Fremdsprachen durchführen (Englisch-Lehrkräfte korrigieren dann Abitur-Klausuren).}{\textbf{Recommendation:} Conduct the KLE in the 2nd semester for 2nd/3rd foreign languages (English teachers will then be correcting Abitur exams).}

\bisection{Prüfungsformat KLE (20 Minuten)}{KLE Examination Format (20 minutes)}

\begin{center}
\begin{tabularx}{\textwidth}{|l|l|X|c|}
\hline
\rowcolor{spanischrot}\color{white}\textbf{Teil} & \color{white}\textbf{Dauer} & \color{white}\textbf{Inhalt / \textit{Content}} & \color{white}\textbf{Bewertet?} \\
\hline
\textbf{Teil 1} & ca. 2 Min. & \textbf{Interview / Aufwärmphase}\newline {\color{englishgray}\textit{Warm-up phase}}\newline Je 2 Fragen pro Kandidat/in, allgemeine Alltagsthemen & \xmark \\
\hline
\textbf{Teil 2} & ca. 4 Min.\newline pro Kand. & \textbf{Monologisches Sprechen}\newline {\color{englishgray}\textit{Monological speaking}}\newline Impulsgesteuertes Sprechen (Bild, Cartoon, Grafik, These)\newline Unterschiedliche Impulse zum gleichen Thema & \cmark \\
\hline
\textbf{Teil 3} & ca. 5 Min.\newline pro Kand. & \textbf{Dialogisches Sprechen}\newline {\color{englishgray}\textit{Dialogical speaking}}\newline Situativ eingebettete Diskussion\newline Thematisch anknüpfend an Teil 2 & \cmark \\
\hline
\textbf{Bewertung} & ca. 10 Min. & Jury berät sich, füllt Bewertungsbögen aus\newline {\color{englishgray}\textit{Jury deliberates, fills out assessment forms}} & --- \\
\hline
\end{tabularx}
\end{center}

\vspace{0.3cm}
\bipar{\textbf{Gruppengröße:} Tandem (2 Personen), max. Dreiergruppe \quad | \quad \textbf{Hilfsmittel:} Keine}{\textbf{Group size:} Pair (2 persons), max. group of three \quad | \quad \textbf{Aids:} None}

\bisection{Zielniveaus nach GER}{Target Levels (CEFR)}

\begin{center}
\begin{tabular}{|l|c|}
\hline
\rowcolor{spanischrot}\color{white}\textbf{Kursart / Phase} & \color{white}\textbf{Zielniveau} \\
\hline
Einführungsphase, fortgeführte FS (ab Kl. 7) & B1 \\
\hline
Einführungsphase, neu begonnene FS (ab Kl. 11) & A2 \\
\hline
Qualifikationsphase, fortgeführte FS & \textbf{B2} \\
\hline
Qualifikationsphase, neu begonnene FS & \textbf{B1(+)} \\
\hline
\end{tabular}
\end{center}

\bisection{Bewertungskriterien}{Assessment Criteria}

\begin{center}
\begin{tikzpicture}
% Gesamtnote Formel
\node[draw=spanischrot, line width=2pt, rounded corners, fill=lightgray, minimum width=12cm, minimum height=1cm] (formula) {
\Large\textbf{GESAMTNOTE = (Sprachliche Leistung $\times$ 60\%) + (Inhaltliche Leistung $\times$ 40\%)}
};
\end{tikzpicture}
\end{center}

\begin{multicols}{2}
\begin{infobox}[Sprachliche Leistung (60\%) / \textit{Linguistic Performance}]
\begin{itemize}[leftmargin=*]
\item Flüssigkeit und Interaktion\\ {\color{englishgray}\textit{Fluency and interaction}}
\item Spektrum sprachlicher Mittel\\ {\color{englishgray}\textit{Range of linguistic means}}
\item Korrektheit (Grammatik, Aussprache)\\ {\color{englishgray}\textit{Accuracy (grammar, pronunciation)}}
\item Kommunikative Strategien\\ {\color{englishgray}\textit{Communicative strategies}}
\item Diskursive Fähigkeiten\\ {\color{englishgray}\textit{Discourse skills}}
\end{itemize}
\end{infobox}

\columnbreak

\begin{infobox}[Inhaltliche Leistung (40\%) / \textit{Content Performance}]
\begin{itemize}[leftmargin=*]
\item Aufgabenerfüllung\\ {\color{englishgray}\textit{Task completion}}
\item Themenbezug\\ {\color{englishgray}\textit{Topic relevance}}
\item Argumentation / Kohärenz\\ {\color{englishgray}\textit{Argumentation / coherence}}
\end{itemize}
\vspace{1.5cm}
\end{infobox}
\end{multicols}

\bisubsection{Tolerierte Besonderheiten mündlicher Kommunikation}{Tolerated Features of Oral Communication}

\bipar{Folgende Aspekte sind \textbf{zu tolerieren} und werden nicht negativ bewertet:}{The following aspects \textbf{must be tolerated} and are not assessed negatively:}

\begin{multicols}{2}
\begin{itemize}
\item Unvollständige Sätze\\ {\color{englishgray}\textit{Incomplete sentences}}
\item Brüche in der Konstruktion\\ {\color{englishgray}\textit{Breaks in construction}}
\item Weniger komplexe Strukturen\\ {\color{englishgray}\textit{Less complex structures}}
\item Geringerer Abstraktionsgrad\\ {\color{englishgray}\textit{Lower level of abstraction}}
\item Wortwiederholungen\\ {\color{englishgray}\textit{Word repetitions}}
\item Natürliche Sprechpausen\\ {\color{englishgray}\textit{Natural speech pauses}}
\item Füllwörter (bueno, pues, entonces...)\\ {\color{englishgray}\textit{Filler words}}
\end{itemize}
\end{multicols}

\bisection{Bewertungsbögen (RECHTLICH VERBINDLICH!)}{Assessment Forms (LEGALLY MANDATORY!)}

\begin{kritischbox}
\textbf{ACHTUNG: Die Excel-Bewertungsbögen (Anlage II)\footnote{Anlage II = Anhang II der Handreichung / \textit{Appendix II of the guidelines}} sind RECHTLICH VERBINDLICH!}\\[0.3cm]
{\color{englishgray}\textit{WARNING: The Excel assessment forms (Appendix II) are LEGALLY MANDATORY!}}\\[0.5cm]

\textbf{Rechtsgrundlage:}\\
Handreichung zur KLE\footnote{KLE = Komplexe Leistungsermittlung / \textit{Complex Performance Assessment}} Sprechen, S. 13:\\
\glqq Die \textbf{verbindlich zu nutzenden} Bewertungsbögen für Gesprächsleitungen und Protokollanten (Anlage II) sind im SIP-Datentauschportal\footnote{SIP = Schulinformations- und Planungssystem M-V / \textit{School Information and Planning System M-V}} unter `Sprechleistung\_Moderne Fremdsprachen' abrufbar.\grqq\\[0.3cm]
{\color{englishgray}\textit{Legal basis: Guidelines for KLE Speaking, p. 13: ``The \textbf{mandatory} assessment forms [...] are available in the SIP data exchange portal.''}}\\[0.5cm]

\textbf{Konsequenz bei Nichtverwendung:}\\
$\bullet$ Prüfung ist \textbf{formal anfechtbar}\\
$\bullet$ Noten können bei Widerspruch \textbf{aufgehoben} werden\\
$\bullet$ \textbf{Eigene Bewertungsbögen sind NICHT zulässig!}\\[0.3cm]
{\color{englishgray}\textit{Consequences of non-use: Examination is formally contestable; grades can be annulled upon appeal; custom assessment forms are NOT permitted!}}
\end{kritischbox}

\vspace{0.5cm}

\bisubsection{Wie erhalte ich die Bewertungsbögen?}{How do I obtain the assessment forms?}

\bipar{Da die Excel-Dateien nur über schulinterne Portale verfügbar sind, gibt es folgende Beschaffungswege:}{Since the Excel files are only available through school-internal portals, there are the following ways to obtain them:}

\begin{center}
\begin{tabularx}{\textwidth}{|c|l|X|}
\hline
\rowcolor{spanischrot}\color{white}\textbf{Nr.} & \color{white}\textbf{Beschaffungsweg} & \color{white}\textbf{Details / \textit{Details}} \\
\hline
\textbf{1} & \textbf{SIP-Datentauschportal} & Zugang über Schulserver / Schulleitung anfragen\newline {\color{englishgray}\textit{Access via school server / ask school administration}}\newline Pfad: \texttt{Sprechleistung\_Moderne Fremdsprachen} \\
\hline
\textbf{2} & \textbf{itslearning-Kurs} & Kurs: \glqq Sprechen testen in den Fremdsprachen SEK II\grqq\newline {\color{englishgray}\textit{Course: ``Testing Speaking in Foreign Languages Upper Secondary''}}\newline Enthält: Anlagen I, III, IV + ggf. Bewertungsbögen \\
\hline
\textbf{3} & \textbf{Schulleitung} & Oberstufenkoordinator/in oder Schulleiter/in\newline {\color{englishgray}\textit{Upper school coordinator or principal}}\newline Haben Zugang zu allen Portalen \\
\hline
\textbf{4} & \textbf{Fachkonferenzleitung} & FK-Leitung\footnote{FK = Fachkonferenz / \textit{Subject Conference}} Fremdsprachen\newline {\color{englishgray}\textit{Subject conference leader for foreign languages}}\newline Sollte Materialien für alle Kolleg/innen haben \\
\hline
\textbf{5} & \textbf{IQ M-V\footnote{IQ M-V = Institut für Qualitätsentwicklung Mecklenburg-Vorpommern / \textit{Institute for Quality Development M-V}} direkt} & E-Mail: \texttt{iq@bm.mv-regierung.de}\newline {\color{englishgray}\textit{Direct request to the Institute}}\newline Freundlich anfragen -- in der Regel hilfsbereit! \\
\hline
\textbf{6} & \textbf{Fachberater Spanisch} & Über Schulamt erfragen\newline {\color{englishgray}\textit{Inquire through school authority}}\newline Kennt alle offiziellen Bezugsquellen \\
\hline
\end{tabularx}
\end{center}

\vspace{0.3cm}

\begin{warnbox}
\centering
\textbf{EMPFEHLUNG / \textit{RECOMMENDATION}}\\[0.3cm]
Beginnen Sie mit \textbf{Option 3 oder 4} (Schulleitung / Fachkonferenz).\\
Falls erfolglos: \textbf{Option 5} (IQ M-V direkt anschreiben).\\[0.2cm]
{\color{englishgray}\textit{Start with Option 3 or 4 (school administration / subject conference).\\If unsuccessful: Option 5 (contact IQ M-V directly).}}
\end{warnbox}

\vspace{0.5cm}

\bisubsubsection{Struktur der Bewertungsbögen}{Structure of Assessment Forms}

\bipar{Jede Excel-Arbeitsmappe enthält:}{Each Excel workbook contains:}

\begin{center}
\begin{tabularx}{\textwidth}{|l|X|}
\hline
\rowcolor{spanischrot}\color{white}\textbf{Bogen} & \color{white}\textbf{Inhalt / \textit{Content}} \\
\hline
\textbf{Bewertungsbogen Gesprächsleitung} & Kriteriengestützte Bewertung durch die prüfende Lehrkraft während der Prüfung \newline {\color{englishgray}\textit{Criteria-based assessment by the examining teacher during the exam}} \\
\hline
\textbf{Bewertungsbogen I (Protokollant)} & Monologisches Sprechen -- analytische Bewertungsskala \newline {\color{englishgray}\textit{Monological speaking -- analytical assessment scale}} \\
\hline
\textbf{Bewertungsbogen II (Protokollant)} & Dialogisches Sprechen -- analytische Bewertungsskala \newline {\color{englishgray}\textit{Dialogical speaking -- analytical assessment scale}} \\
\hline
\end{tabularx}
\end{center}

\bisubsubsection{Bewertungskategorien (Beispiel B2)}{Assessment Categories (Example B2)}

\begin{center}
\begin{tabularx}{\textwidth}{|l|X|c|}
\hline
\rowcolor{spanischrot}\color{white}\textbf{Kategorie} & \color{white}\textbf{Beschreibung / \textit{Description}} & \color{white}\textbf{Punkte} \\
\hline
\multicolumn{3}{|l|}{\cellcolor{lightgray}\textbf{Sprachliche Leistung (60\%)}} \\
\hline
Flüssigkeit & Redefluss, Tempo, natürliche Pausen \newline {\color{englishgray}\textit{Flow of speech, pace, natural pauses}} & 0--5 \\
\hline
Interaktion & Eingehen auf Partner, Gesprächsführung \newline {\color{englishgray}\textit{Responding to partner, conversation management}} & 0--5 \\
\hline
Spektrum & Wortschatz, grammatische Strukturen \newline {\color{englishgray}\textit{Vocabulary, grammatical structures}} & 0--5 \\
\hline
Korrektheit & Grammatik, Aussprache, Intonation \newline {\color{englishgray}\textit{Grammar, pronunciation, intonation}} & 0--5 \\
\hline
\multicolumn{3}{|l|}{\cellcolor{lightgray}\textbf{Inhaltliche Leistung (40\%)}} \\
\hline
Aufgabenerfüllung & Vollständigkeit, Relevanz \newline {\color{englishgray}\textit{Completeness, relevance}} & 0--5 \\
\hline
Themenbezug & Sachkenntnis, Bezug zum Thema \newline {\color{englishgray}\textit{Subject knowledge, topic relevance}} & 0--5 \\
\hline
Argumentation & Kohärenz, Struktur, Begründungen \newline {\color{englishgray}\textit{Coherence, structure, reasoning}} & 0--5 \\
\hline
\end{tabularx}
\end{center}

\bisection{Organisation}{Organization}

\bisubsection{Prüfungsausschuss}{Examination Committee}

\begin{center}
\begin{tabularx}{\textwidth}{|l|X|}
\hline
\rowcolor{spanischrot}\color{white}\textbf{Rolle / \textit{Role}} & \color{white}\textbf{Aufgaben / \textit{Tasks}} \\
\hline
\textbf{Gesprächsleitung} \newline {\color{englishgray}\textit{Discussion Leader}} &
$\bullet$ Direkter Blickkontakt mit Kandidaten \newline
$\bullet$ Zeit kontrollieren \newline
$\bullet$ Ggf. Zusatzimpuls geben (bei Black-out) \newline
$\bullet$ Mienenspiel kontrollieren \newline
$\bullet$ \textbf{Keine inhaltliche/sprachliche Hilfe!} \newline
{\color{englishgray}\textit{No content or language assistance!}} \\
\hline
\textbf{Protokollant/in} \newline {\color{englishgray}\textit{Note-taker}} &
$\bullet$ Seitliche Position \newline
$\bullet$ Passive Beobachtung \newline
$\bullet$ Bewertungsbogen ausfüllen \newline
$\bullet$ Andere Perspektive als Gesprächsleitung \\
\hline
\end{tabularx}
\end{center}

\bisubsection{Tandem-Zusammenstellung}{Pair Formation}

\bipar{\textbf{Optionen:}}{Options:}
\begin{enumerate}
\item Lehrkraft bestimmt (leistungshomogen oder -heterogen)\\ {\color{englishgray}\textit{Teacher determines (homogeneous or heterogeneous performance levels)}}
\item SuS wählen selbst (Abstimmung mit Lehrkraft erforderlich)\\ {\color{englishgray}\textit{Students choose themselves (coordination with teacher required)}}
\item Gemeinsame Festlegung Lehrkraft + SuS\\ {\color{englishgray}\textit{Joint determination by teacher and students}}
\item Losverfahren\\ {\color{englishgray}\textit{Random draw}}
\end{enumerate}

\begin{warnbox}
\textbf{Bekanntgabe:} Spätestens \textbf{1 Tag vor} der Prüfung\\
{\color{englishgray}\textit{Announcement: At least \textbf{1 day before} the examination}}
\end{warnbox}

\bisubsection{Zeitplanung (Beispiel)}{Time Schedule (Example)}

\bipar{Beispiel für 22 Kandidaten, 1 Jury (2 Fachlehrkräfte):}{Example for 22 candidates, 1 jury (2 subject teachers):}

\begin{center}
\begin{tabular}{|l|l|}
\hline
\rowcolor{spanischrot}\color{white}\textbf{Zeit} & \color{white}\textbf{Aktivität / \textit{Activity}} \\
\hline
08:00 -- 08:30 & Tandem 1 \\
08:30 -- 09:00 & Tandem 2 \\
09:00 -- 09:30 & Tandem 3 \\
09:30 -- 10:00 & Tandem 4 \\
\hline
10:00 -- 10:15 & \textbf{Pause / \textit{Break}} \\
\hline
10:15 -- 10:45 & Tandem 5 \\
... & ... \\
\hline
$\approx$ 14:15 & Ende (mit 30 Min. Mittagspause) \\
\hline
\end{tabular}
\end{center}

\bipar{\textbf{Bei Dreiergruppen:} 30 Min. Prüfung + 10 Min. Bewertung = 40 Min. pro Gruppe}{\textbf{For groups of three:} 30 min. exam + 10 min. assessment = 40 min. per group}

\bisubsection{Sonderfall: Audiographie}{Special Case: Audio Recording}

\bipar{In begründeten Ausnahmefällen (keine 2. Fachlehrkraft verfügbar) kann eine Audiographie eingesetzt werden:}{In justified exceptional cases (no second subject teacher available), audio recording may be used:}

\begin{itemize}
\item Nur \textbf{einmal} anhören / {\color{englishgray}\textit{Listen only \textbf{once}}}
\item Korrekturzeitraum: max. 3 Wochen / {\color{englishgray}\textit{Correction period: max. 3 weeks}}
\item Danach \textbf{sofort vernichten} / {\color{englishgray}\textit{Then \textbf{destroy immediately}}}
\end{itemize}

% ============================================================================
% TEIL B: ABITUR-PAARPRÜFUNG
% ============================================================================
\newpage
\part*{\color{spanischrot}Teil B: Abitur-Paarprüfung}
\addcontentsline{toc}{part}{Teil B: Abitur-Paarprüfung}
{\large\color{englishgray}\textit{Part B: Abitur Pair Examination}}

\bisection{Rechtliche Grundlagen}{Legal Framework}

\begin{center}
\begin{tabularx}{\textwidth}{|l|X|}
\hline
\rowcolor{spanischrot}\color{white}\textbf{Paragraph} & \color{white}\textbf{Regelung / \textit{Regulation}} \\
\hline
\textbf{§ 25 Abs. 8 APVO} & Optionaler mündlicher Prüfungsteil zur \textbf{verkürzten} schriftlichen Abiturprüfung \newline {\color{englishgray}\textit{Optional oral examination component for \textbf{shortened} written Abitur exam}} \\
\hline
\textbf{§ 38 Abs. 6 APVO} & Mündliche Prüfung als Paarprüfung möglich (4./5. Prüfungsfach) \newline {\color{englishgray}\textit{Oral exam as pair examination possible (4th/5th exam subject)}} \\
\hline
\textbf{§ 28 Abs. 3 APVO} & Fachprüfungsausschuss mit nur 2 Mitgliedern genehmigt \newline {\color{englishgray}\textit{Examination committee with only 2 members approved}} \\
\hline
\end{tabularx}
\end{center}

\vspace{0.3cm}
\bipar{\textbf{Gültig seit:} 11.12.2025 (Handreichung, VV vom 03.12.2025)}{\textbf{Valid since:} December 11, 2025 (Guidelines, Administrative Regulation dated December 3, 2025)}

\begin{warnbox}
\textbf{Voraussetzung:} Schule hat Format eingeführt (Schulversuch ODER Schulkonferenzbeschluss)\\
{\color{englishgray}\textit{Prerequisite: School has introduced the format (pilot project OR school conference decision)}}
\end{warnbox}

\bisection{Zwei Varianten}{Two Variants}

\begin{center}
\begin{tabularx}{\textwidth}{|X|X|}
\hline
\rowcolor{spanischrot}\color{white}\textbf{Variante 1: § 25 Abs. 8} & \color{white}\textbf{Variante 2: § 38 Abs. 6} \\
\hline
Mündlicher Prüfungsteil + \textbf{verkürzte} schriftliche Abiturprüfung \newline {\color{englishgray}\textit{Oral exam component + \textbf{shortened} written Abitur exam}} &
Mündliche Prüfung als 4./5. Prüfungsfach \newline {\color{englishgray}\textit{Oral exam as 4th/5th exam subject}} \\
\hline
Auf Antrag der SuS, Antragsfrist schulintern \newline {\color{englishgray}\textit{Upon student application, deadline set by school}} &
Auf Antrag der SuS, Frist vom Prüfungsvorsitzenden bestimmt \newline {\color{englishgray}\textit{Upon student application, deadline set by examination chair}} \\
\hline
\end{tabularx}
\end{center}

\bisection{Prüfungsformat (max. 30 Minuten)}{Examination Format (max. 30 minutes)}

\begin{center}
\begin{tabularx}{\textwidth}{|l|l|X|c|}
\hline
\rowcolor{spanischrot}\color{white}\textbf{Teil} & \color{white}\textbf{Dauer} & \color{white}\textbf{Inhalt / \textit{Content}} & \color{white}\textbf{Bewertet?} \\
\hline
\textbf{Teil 1} & ca. 2 Min. & \textbf{Interview / Aufwärmphase}\newline {\color{englishgray}\textit{Warm-up phase}}\newline Je 2 Fragen pro Prüfling, allgemeine Alltagsthemen & \xmark \\
\hline
\textbf{Teil 2} & ca. 4 Min.\newline pro Prüfl. & \textbf{Monologisches Sprechen}\newline {\color{englishgray}\textit{Monological speaking}}\newline Unterschiedliche Impulse zum \textbf{gleichen Thema}\newline Bezug zu Rahmenplan und Vorabhinweisen & \cmark \\
\hline
\textbf{Teil 3} & ca. 5 Min.\newline pro Prüfl. & \textbf{Dialogisches Sprechen}\newline {\color{englishgray}\textit{Dialogical speaking}}\newline \textbf{ANDERER Themenschwerpunkt als Teil 2!}\newline Situativ eingebettete Diskussion & \cmark \\
\hline
\end{tabularx}
\end{center}

\begin{warnbox}
\textbf{WICHTIG:} Teil 2 und Teil 3 haben \textbf{unterschiedliche Themenschwerpunkte}!\\
{\color{englishgray}\textit{IMPORTANT: Parts 2 and 3 have \textbf{different topic focuses}!}}\\[0.2cm]
(Dies ist der Hauptunterschied zur KLE, wo beide Teile das gleiche Thema behandeln.)\\
{\color{englishgray}\textit{(This is the main difference from KLE, where both parts cover the same topic.)}}
\end{warnbox}

\bisection{Wesentliche Unterschiede zur KLE}{Key Differences from KLE}

\begin{center}
\begin{tabularx}{\textwidth}{|l|X|X|}
\hline
\rowcolor{spanischrot}\color{white}\textbf{Aspekt} & \color{white}\textbf{Komplexe LE (KLE)} & \color{white}\textbf{Abitur-Paarprüfung} \\
\hline
\textbf{Zeitpunkt} & Q1, Q2 oder Q3 & Nach schriftlichem Abitur \\
\hline
\textbf{Dauer} & 20 Minuten & max. 30 Minuten \\
\hline
\textbf{Pflicht/Optional} & \textbf{PFLICHT} (1x) & Optional (Antrag) \\
\hline
\textbf{Teil 2 \& 3 Thema} & Gleiches Thema & \textbf{Unterschiedliche} Themen \\
\hline
\textbf{Aufgabenerstellung} & Fachlehrkraft & Fachlehrkraft + Vorabhinweise \\
\hline
\textbf{Protokoll} & Bewertungsbögen & Bewertungsbögen + \textbf{Erwartungshorizont} \\
\hline
\textbf{Bekanntgabe Tandem} & 1 Tag vorher & \textbf{2 Werktage} vorher \\
\hline
\textbf{Rechtsgrundlage} & § 22 Abs. 5 APVO & § 25 Abs. 8 / § 38 Abs. 6 APVO \\
\hline
\end{tabularx}
\end{center}

\bisection{Aufgabenerstellung Abitur}{Abitur Task Design}

\bipar{\textbf{Zusätzliche Anforderungen:}}{Additional requirements:}

\begin{itemize}
\item Vorabhinweise des jeweiligen Prüfungsjahres beachten\\ {\color{englishgray}\textit{Observe preliminary notes for the respective examination year}}
\item Rahmenplan Qualifikationsphase beachten\\ {\color{englishgray}\textit{Follow the curriculum for the qualification phase}}
\item Teil 2 und 3: \textbf{unterschiedliche Themenschwerpunkte}\\ {\color{englishgray}\textit{Parts 2 and 3: \textbf{different topic focuses}}}
\item Erwartungshorizont erstellen\\ {\color{englishgray}\textit{Create expectations/marking scheme}}
\item AFB I--III berücksichtigen\\ {\color{englishgray}\textit{Consider requirement levels I--III}}
\end{itemize}

\bisection{Bekanntgabe der Ergebnisse}{Announcement of Results}

\bipar{Gemäß § 44 Abs. 4 APVO M-V:}{According to § 44 Para. 4 APVO M-V:}

\begin{itemize}
\item Am Ende des halben oder ganzen Prüfungstages\\ {\color{englishgray}\textit{At the end of the half or full examination day}}
\item Durch Vorsitzenden der Prüfungskommission\\ {\color{englishgray}\textit{By the chair of the examination committee}}
\item Auf Verlangen: Mündliche Erläuterung der Bewertung\\ {\color{englishgray}\textit{Upon request: Oral explanation of the assessment}}
\item Bei begründeten Einwänden: Eingehen erforderlich\\ {\color{englishgray}\textit{In case of justified objections: Response required}}
\item Keine schriftliche Begründung nötig\\ {\color{englishgray}\textit{No written explanation required}}
\end{itemize}

% ============================================================================
% TEIL C: CHECKLISTEN
% ============================================================================
\newpage
\part*{\color{spanischrot}Teil C: Checklisten}
\addcontentsline{toc}{part}{Teil C: Checklisten}
{\large\color{englishgray}\textit{Part C: Checklists}}

\bisection{Checkliste: KLE vorbereiten}{Checklist: Preparing KLE}

\bisubsection{Planung (mehrere Wochen vorher)}{Planning (several weeks in advance)}

\begin{itemize}[label=$\square$]
\item Termin mit Schulleitung/Oberstufenkoordination abstimmen\\ {\color{englishgray}\textit{Coordinate date with school administration/upper school coordinator}}
\item Zweitprüfer/in (Protokollant) organisieren\\ {\color{englishgray}\textit{Organize second examiner (note-taker)}}
\item Termin in Klausurplan integrieren\\ {\color{englishgray}\textit{Integrate date into exam schedule}}
\item Kollegium informieren (Lehrerkonferenz)\\ {\color{englishgray}\textit{Inform staff (teachers' conference)}}
\end{itemize}

\bisubsection{Aufgabenerstellung (2--3 Wochen vorher)}{Task Design (2--3 weeks in advance)}

\begin{itemize}[label=$\square$]
\item Thema rahmenplankonform wählen\\ {\color{englishgray}\textit{Choose topic in accordance with curriculum}}
\item Impulse für Teil 2 erstellen (Bilder, Cartoons, Grafiken...)\\ {\color{englishgray}\textit{Create prompts for Part 2 (images, cartoons, graphs...)}}
\item Aufgabe für Teil 3 erstellen (situativ eingebettet)\\ {\color{englishgray}\textit{Create task for Part 3 (situationally embedded)}}
\item Zusatzimpulse für Notfälle vorbereiten\\ {\color{englishgray}\textit{Prepare backup prompts for emergencies}}
\item Vokabular und Strukturen auf Bekanntes prüfen\\ {\color{englishgray}\textit{Check vocabulary and structures are familiar}}
\item Bewertungsbögen (A2/B1/B2) vorbereiten\\ {\color{englishgray}\textit{Prepare assessment forms (A2/B1/B2)}}
\end{itemize}

\bisubsection{Organisation (1 Woche vorher)}{Organization (1 week in advance)}

\begin{itemize}[label=$\square$]
\item Tandems zusammenstellen (pädagogische Kriterien)\\ {\color{englishgray}\textit{Form pairs (pedagogical criteria)}}
\item Zeitplan erstellen (30 Min. pro Tandem)\\ {\color{englishgray}\textit{Create schedule (30 min. per pair)}}
\item Prüfungsraum und Aufenthaltsraum organisieren\\ {\color{englishgray}\textit{Organize examination room and waiting room}}
\end{itemize}

\bisubsection{Kurz vorher}{Shortly before}

\begin{itemize}[label=$\square$]
\item Tandems bekannt geben (1 Tag vorher)\\ {\color{englishgray}\textit{Announce pairs (1 day before)}}
\item SuS über Ablauf und Bewertung informieren\\ {\color{englishgray}\textit{Inform students about procedure and assessment}}
\item Regieanweisungen lesen\\ {\color{englishgray}\textit{Read stage directions}}
\end{itemize}

\bisubsection{Am Prüfungstag}{On the examination day}

\begin{itemize}[label=$\square$]
\item Räume vorbereiten\\ {\color{englishgray}\textit{Prepare rooms}}
\item Bewertungsbögen bereithalten\\ {\color{englishgray}\textit{Have assessment forms ready}}
\item Timer/Uhr bereithalten\\ {\color{englishgray}\textit{Have timer/clock ready}}
\end{itemize}

\bisection{Checkliste: Abitur-Paarprüfung}{Checklist: Abitur Pair Examination}

\bisubsection{Voraussetzungen}{Prerequisites}

\begin{itemize}[label=$\square$]
\item Schule hat Format eingeführt (Schulkonferenzbeschluss)\\ {\color{englishgray}\textit{School has introduced format (school conference decision)}}
\item Prüflingszahlen an IQ M-V gemeldet\\ {\color{englishgray}\textit{Number of candidates reported to IQ M-V}}
\end{itemize}

\bisubsection{Aufgabenerstellung}{Task Design}

\begin{itemize}[label=$\square$]
\item Vorabhinweise des Prüfungsjahres beachtet\\ {\color{englishgray}\textit{Preliminary notes for examination year observed}}
\item Teil 2 und 3: \textbf{unterschiedliche Themenschwerpunkte}\\ {\color{englishgray}\textit{Parts 2 and 3: \textbf{different topic focuses}}}
\item Erwartungshorizont erstellt\\ {\color{englishgray}\textit{Expectations/marking scheme created}}
\item AFB I--III berücksichtigt\\ {\color{englishgray}\textit{Requirement levels I--III considered}}
\end{itemize}

\bisubsection{Organisation}{Organization}

\begin{itemize}[label=$\square$]
\item Tandems \textbf{2 Werktage} vorher bekannt gegeben\\ {\color{englishgray}\textit{Pairs announced \textbf{2 working days} in advance}}
\item Bewertungsbögen B1 oder B2 vorbereitet\\ {\color{englishgray}\textit{Assessment forms B1 or B2 prepared}}
\item Protokollierung vorbereitet\\ {\color{englishgray}\textit{Documentation prepared}}
\end{itemize}

% ============================================================================
% TEIL D: REGIEANWEISUNGEN
% ============================================================================
\newpage
\part*{\color{spanischrot}Teil D: Regieanweisungen}
\addcontentsline{toc}{part}{Teil D: Regieanweisungen}
{\large\color{englishgray}\textit{Part D: Stage Directions}}

\bipar{Die folgenden Regieanweisungen sind für Spanisch angepasst. Das englische Original findet sich in Anlage III der Handreichung.}{The following stage directions are adapted for Spanish. The English original can be found in Appendix III of the guidelines.}

\bisection{Teil 1: Interview (Aufwärmphase)}{Part 1: Interview (Warm-up Phase)}

\begin{tcolorbox}[colback=lightgray, colframe=spanischrot, title=Gesprächsleitung / \textit{Discussion Leader}]
\textbf{Español:}\\
„En la primera parte de este examen, voy a haceros algunas preguntas.\\
Las preguntas son sobre [Thema].\\
Empezamos."\\[0.3cm]
\textit{[Alternierend A-B/B-A/A-B Fragen stellen]}\\[0.3cm]
{\color{englishgray}\textbf{English:}\\
``In the first part of this test, I'm going to ask each of you some questions.\\
The questions are about [topic].\\
Let's begin.''}
\end{tcolorbox}

\bisection{Teil 2: Monologisches Sprechen}{Part 2: Monological Speaking}

\begin{tcolorbox}[colback=lightgray, colframe=spanischrot, title=Gesprächsleitung / \textit{Discussion Leader}]
\textbf{Español:}\\
„Ahora me gustaría que cada uno/a hable sobre [Thema] (imágenes, diagramas...).\\
Tenéis aproximadamente 4 minutos cada uno/a.\\[0.2cm]
Candidato/a A [Name]: Aquí está tu material.\\
¿Quieres que repita las instrucciones?"\\[0.2cm]
\textit{[Falls ja: Anweisungen wiederholen]}\\
\textit{[Zeit zum Betrachten des Materials geben, falls nötig]}\\[0.2cm]
„Por favor, lee ahora la tarea para el habla monológica.\\
¿Estás listo/a para empezar?"\\[0.2cm]
\textit{[Bei Bedarf: Zusatzimpuls für Notfall stellen]}\\
\textit{[Nach 4 Min.: Wechsel zu Kandidat B]}\\[0.2cm]
„Candidato/a B [Name]: Aquí está tu material.\\
¿Quieres que repita las instrucciones?"\\[0.3cm]
{\color{englishgray}\textbf{English:}\\
``Now I'd like each of you to talk on your own about [topic] (pictures, diagrams...).\\
You have about 4 minutes each.\\
Candidate A [name]: Here is your material.\\
Shall I repeat the instructions?''}
\end{tcolorbox}

\bisection{Teil 3: Dialogisches Sprechen}{Part 3: Dialogical Speaking}

\begin{tcolorbox}[colback=lightgray, colframe=spanischrot, title=Gesprächsleitung / \textit{Discussion Leader}]
\textbf{Español:}\\
„Ahora me gustaría que habléis entre vosotros/as.\\
Por favor, coged vuestro material y leed ahora la tarea para el habla dialógica.\\
¿Queréis que repita las instrucciones?\\
¿Estáis listos/as para empezar?"\\[0.3cm]
\textit{[Falls niemand beginnt:]}\\
„Candidato/a A/B [Name], por favor, empieza."\\[0.2cm]
\textit{[Bei überproportionalem Redeanteil:]}\\
„Gracias, candidato/a A, ¿y tú, candidato/a B?"\\[0.2cm]
\textit{[Bei längeren Redepausen (20--30 Sek.):]}\\
„¿Y tú, candidato/a B?... Vale. Candidato/a A, ¿tienes más ideas?"\\[0.2cm]
\textit{[Nach Ablauf der Zeit:]}\\
„Gracias. Esto es el final de vuestro examen oral.\\
Por favor, volved inmediatamente a vuestras clases."\\[0.3cm]
{\color{englishgray}\textbf{English:}\\
``Now I'd like you to talk to each other.\\
Please take your material and read now the task for dialogical speaking.\\
Shall I repeat the instructions?\\
Are you ready to begin?''}
\end{tcolorbox}

% ============================================================================
% TEIL E: FORMULARE
% ============================================================================
\newpage
\part*{\color{spanischrot}Teil E: Formulare und Vorlagen}
\addcontentsline{toc}{part}{Teil E: Formulare und Vorlagen}
{\large\color{englishgray}\textit{Part E: Forms and Templates}}

\bisection{Ergebnismitteilung}{Result Notification}

\bipar{Die folgende Vorlage (Anlage IV der Handreichung) wird nach der KLE an die SuS ausgegeben:}{The following template (Appendix IV of the guidelines) is given to students after the KLE:}

\begin{tcolorbox}[colback=white, colframe=spanischrot, boxrule=1.5pt]
\begin{center}
\rule{\textwidth}{0.5pt}\\[0.3cm]
{\small Name der Schule / \textit{School name}}\\[0.5cm]

\underline{\hspace{8cm}} {\small (Name des Schülers / der Schülerin)}\\[0.3cm]
{\color{englishgray}\textit{(Student name)}}\\[0.5cm]

hat in der komplexen Leistungsermittlung im Kompetenzbereich Sprechen\\
{\color{englishgray}\textit{has achieved in the complex assessment of speaking skills}}\\[0.3cm]

am \underline{\hspace{3cm}} (Datum) um \underline{\hspace{2cm}} (Uhrzeit)\\
{\color{englishgray}\textit{on (date) at (time)}}\\[0.3cm]

zum Thema \underline{\hspace{6cm}}\\
{\color{englishgray}\textit{on the topic of}}\\[0.3cm]

\underline{\hspace{2cm}} / \underline{\hspace{2cm}} (Notenpunkte/Note) erreicht.\\
{\color{englishgray}\textit{(grade points/grade) achieved.}}\\[0.5cm]

\rule{\textwidth}{0.5pt}\\[0.3cm]
{\small Ort, Datum \hfill Unterschrift der unterrichtenden Fachlehrkraft}\\
{\small\color{englishgray}\textit{Place, date \hfill Signature of subject teacher}}\\[0.5cm]

\rule{\textwidth}{0.5pt}\\[0.3cm]
{\small Ort, Datum \hfill Unterschrift volljährige*r Schüler*in bzw. der Erziehungsberechtigten}\\
{\small\color{englishgray}\textit{Place, date \hfill Signature of adult student or guardian}}
\end{center}
\end{tcolorbox}

% ============================================================================
% TEIL F: QUELLEN
% ============================================================================
\newpage
\part*{\color{spanischrot}Teil F: Quellen und Rechtsgrundlagen}
\addcontentsline{toc}{part}{Teil F: Quellen und Rechtsgrundlagen}
{\large\color{englishgray}\textit{Part F: Sources and Legal Framework}}

\bisection{Handreichungen}{Guidelines}

\begin{center}
\begin{tabularx}{\textwidth}{|X|c|}
\hline
\rowcolor{spanischrot}\color{white}\textbf{Dokument / \textit{Document}} & \color{white}\textbf{Gültig ab} \\
\hline
Handreichung komplexe Leistungsermittlung Sprechen\newline {\color{englishgray}\textit{Guidelines for Complex Assessment of Speaking Skills}} & 01.08.2020 \\
\hline
Handreichung Paarprüfung Abitur\newline {\color{englishgray}\textit{Guidelines for Abitur Pair Examination}} & 11.12.2025 \\
\hline
\end{tabularx}
\end{center}

\bisection{Verwaltungsvorschriften}{Administrative Regulations}

\begin{center}
\begin{tabularx}{\textwidth}{|X|c|}
\hline
\rowcolor{spanischrot}\color{white}\textbf{Dokument / \textit{Document}} & \color{white}\textbf{Datum} \\
\hline
VV Durchführung Sprechprüfung (Mittl.bl. BM M-V Nr. 5/2020, S. 290)\newline {\color{englishgray}\textit{AR Implementation Speaking Examination}} & 24.07.2020 \\
\hline
VV Paarprüfung Abitur (Mittl.bl. BM M-V Nr. 14/2025, S. 239)\newline {\color{englishgray}\textit{AR Abitur Pair Examination}} & 03.12.2025 \\
\hline
\end{tabularx}
\end{center}

\bisection{APVO M-V (Abiturprüfungsverordnung)}{APVO M-V (Abitur Examination Regulations)}

\begin{center}
\begin{tabularx}{\textwidth}{|l|X|}
\hline
\rowcolor{spanischrot}\color{white}\textbf{Paragraph} & \color{white}\textbf{Inhalt / \textit{Content}} \\
\hline
§ 16 Abs. 2 & Anforderungsbereiche I--III / {\color{englishgray}\textit{Requirement levels I--III}} \\
\hline
§ 17 Abs. 3 & Bewertung komplexer Leistungen wie Klausuren / {\color{englishgray}\textit{Grading of complex assessments like written exams}} \\
\hline
§ 21 Abs. 4 & Gewichtung Klausuren/Sonstige (3+: 50/50) / {\color{englishgray}\textit{Weighting exams/other}} \\
\hline
§ 22 Abs. 1 & Gewichtung Klausuren/Sonstige / {\color{englishgray}\textit{Weighting exams/other}} \\
\hline
§ 22 Abs. 2 & Abiturähnliche Klausur Q4 / {\color{englishgray}\textit{Abitur-style exam Q4}} \\
\hline
§ 22 Abs. 4 & Max. 1/3 Klausuren ersetzbar / {\color{englishgray}\textit{Max. 1/3 exams replaceable}} \\
\hline
\textbf{§ 22 Abs. 5} & \textbf{PFLICHT: Mind. 1 KLE Sprechen in Q-Phase} / {\color{englishgray}\textit{MANDATORY: Min. 1 KLE Speaking in Q-phase}} \\
\hline
§ 22 Abs. 6 & Mündliche Überprüfung in Q4 auf Antrag / {\color{englishgray}\textit{Oral exam in Q4 upon request}} \\
\hline
\textbf{§ 25 Abs. 8} & \textbf{Optionaler mündlicher Prüfungsteil Abitur} / {\color{englishgray}\textit{Optional oral exam component Abitur}} \\
\hline
§ 28 Abs. 3 & Fachprüfungsausschuss mit 2 Mitgliedern / {\color{englishgray}\textit{Exam committee with 2 members}} \\
\hline
\textbf{§ 38 Abs. 6} & \textbf{Mündliche Prüfung als Paarprüfung} / {\color{englishgray}\textit{Oral exam as pair examination}} \\
\hline
§ 44 Abs. 4 & Bekanntgabe Prüfungsergebnisse / {\color{englishgray}\textit{Announcement of results}} \\
\hline
§ 55 Abs. 3 & Ausnahme Fachgymnasien / {\color{englishgray}\textit{Exception vocational high schools}} \\
\hline
\end{tabularx}
\end{center}

\bisection{Download-Quellen}{Download Sources}

\begin{infobox}[Materialien / \textit{Materials}]
\textbf{Bewertungsbögen A2/B1/B2:}\\
SIP-Datentauschportal $\rightarrow$ „Sprechleistung\_Moderne Fremdsprachen"\\[0.3cm]
\textbf{Regieanweisungen, Beispielplanungen:}\\
itslearning-Kurs „Sprechen testen in den Fremdsprachen SEK II"\\[0.3cm]
\textbf{Beispielaufgaben Abitur:}\\
itslearning-Kurs „Paarprüfung im Abitur der modernen Fremdsprachen"\\[0.3cm]
\textbf{Rahmenpläne und Handreichungen:}\\
\url{https://www.bildung-mv.de/schueler/schule-und-unterricht/faecher-und-rahmenplaene/}
\end{infobox}

\vfill
\begin{center}
\rule{8cm}{0.5pt}\\[0.3cm]
{\small Erstellt auf Grundlage der offiziellen Dokumente des Ministeriums für Bildung und Kindertagesförderung Mecklenburg-Vorpommern}\\[0.2cm]
{\small\color{englishgray}\textit{Created based on official documents from the Ministry of Education MV}}\\[0.5cm]
{\footnotesize Januar 2026}
\end{center}

\end{document}
